\documentclass[10pt,a4paper]{amsart}
\usepackage{amsmath,amssymb,amsthm,graphicx}
\usepackage{hyperref}
\newcommand{\C}{\mathbb{C}}
\renewcommand{\H}{\mathbb{H}}
\newcommand{\Z}{\mathbb{Z}}
\newcommand{\T}{\mathbb{T}}
\newcommand{\Q}{\mathbb{Q}}
\newcommand{\N}{\mathbb{N}}
\newcommand{\R}{\mathbb{R}}
\newcommand{\slz}{{\text {\rm SL}}_2(\mathbb{Z})}
\newcommand{\re}{\textnormal{Re}}
\newcommand{\im}{\textnormal{Im}}
\newcommand{\Tr}{\operatorname{Tr}}

\newcommand{\vt}[1]{\left\lvert #1 \right\rvert}
\newcommand{\rf}{\mathfrak{r}}
\newcommand{\Ec}{\mathcal{E}}
\newcommand{\Cc}{\mathcal{C}}
\newcommand{\Qc}{\mathcal{Q}}
\newcommand{\Fc}{\mathcal{F}}

\numberwithin{equation}{section}

\newtheorem{Theorem}{Theorem}[section]
\newtheorem*{Theorem*}{Theorem}
\newtheorem{Corollary}[Theorem]{Corollary}
\newtheorem{Lemma}[Theorem]{Lemma}
\newtheorem{Proposition}[Theorem]{Proposition}
{ \theoremstyle{definition}
	\newtheorem{Definition}[Theorem]{Definition}
	\newtheorem{Note}[Theorem]{Note}
	\newtheorem{Example}[Theorem]{Example}
	\newtheorem{Remark}[Theorem]{Remark} }
\usepackage{cleveref}
\crefname{originalreference}{original reference}{original references}
\Crefname{originalreference}{Original reference}{Original references}
\makeatletter
\newlabel{lemma11}{{lemma11}{source}{}{sourcepointer-1}{}}
\newlabel{lemma11@cref}{{[originalreference][0][]lemma11}{[1][0][]source}{}{}{}}
\makeatother
\title{Asymptotics of F\_\{a,c\}(zeta; exp(-z)) as z tends to 0 in D\_theta, for zeta a primitive b-th root of unity}
\author{Males, Joshua}
\date{doi:10.3842/SIGMA.2024.001; CC BY 4.0}
\begin{document}
\maketitle
\setcounter{section}{2}
\setcounter{Theorem}{3}
\noindent\textit{Editorial note.} Author's own Proposition environment, counted statement at source lines 231-236 and its own adjacent proof at source lines 237-273, both read verbatim from sources/sigma-2024-001-a-note-on-the-equidistribution-of-3-colour-p/pinned-original-source.tex. The counted spans are the author's environments as printed; no line was dropped from either span.\par\medskip
\noindent\textit{Editorial note.} Asymmetry of the census record for this unit, recorded so that no reader takes a mispaired block for an independent proof: the census pairs the Proposition labelled Prop: Craig (source lines 276-282) with the proof block at source lines 310-341, but that block is the proof of Theorem 1.1 of this article (its first line names that theorem) and the Prop: Craig statement is quoted from the reference Craig without a proof, so the block belongs to Theorem 1.1 and not to Prop: Craig. No count is taken from Prop: Craig.\par\medskip
\noindent\textit{Editorial note.} The counted proof span contains two citation tokens, the reference to the published digamma identity (Campbell, page 39) and the article's own earlier reference label, both of which are delivered as external references of this card rather than as reproduced text; the local span check reports them and the report is accepted rather than edited away, because the author's proof is delivered whole.\par\medskip
\noindent\textit{Editorial note.} The article is distributed under CC BY 4.0 with the fixed version URL https://sigma-journal.com/2024/001/ and the licence observation is filed next to the source as license-observation.json. The counter seed is taken from the number this Proposition carries in the published article, 2.4 (census version 3.7 field printed): the article's theorem declarations are scoped by section, and three numbered theorem-family environments precede this one inside its section, so the seed is section 2 with the shared theorem counter at 3; the earlier seed built from the pre-3.7 census number 5.1 was wrong and is superseded. The unit was scanned against the published-item list by DOI before build: no published entry carries this DOI. All card fields are ASCII only.\par\medskip
\noindent\textit{Allowed standard prerequisites (editorial).}\begin{itemize}
\item Lerch transcendent and its series representation. Same article, displayed definition at source lines 227-229. Supplies the series Phi(z,s,a) = sum over n of z\textasciicircum{}n/(n+a)\textasciicircum{}s that names the exponential main term of the asymptotic formula.
\item Taylor expansion of the auxiliary function f. Same article, Lemma lemma11 (source lines 169-179), not delivered on this card. Supplies the expansion f(x) = 1/(c x\textasciicircum{}2) + (1/2 - a/c)/x + O(1) at x = m b z, which the proof applies to turn the double sum into the leading logarithm.
\item Digamma identity for roots of unity. Published identity cited in the proof as the reference Campbell, page 39, with the sign correction recorded in the same sentence. Supplies the evaluation of the sum over j of psi(j/b) zeta\textasciicircum{}j as b log(1 - zeta), used when exponentiating the logarithm.
\end{itemize}\medskip
\ifdefined\hypertarget\hypertarget{sourcepointer-1}{}\fi
\noindent\textit{Source reference.} Original reference lemma11: sources/sigma-2024-001-a-note-on-the-equidistribution-of-3-colour-p/pinned-original-source.tex, source lines 169--169. Navigation pointer invoked at proof line 185; the counted span contains the pointer only and the cited result is not reproduced inside it, so the pointer is registered as a non-load-bearing navigation stub in the established form for Stacks material. No bibliography citation is involved.\par\medskip
\begin{Proposition}\label{Prop: F_a,c}
	Let $\zeta$ be a primitive $b$-th root of unity for $b \geq 2$. As $z\to 0$ in $D_\theta$, we have that
	\begin{align*}
		F_{a,c}(\zeta;{\rm e}^{-z}) = (1-\zeta)^{\frac{1}{2} - \frac{a}{c}}\exp\left( -\frac{\zeta \Phi (\zeta,2, 1 )}{ c z} \right) ( 1+ O(|z|) ).
	\end{align*}
\end{Proposition}

\begin{proof}
	First note that
	\begin{align*}
		\log F_{a,c}(\zeta;q) = \sum_{n \geq 0} \log\left(1- \zeta q^{a+cn}\right) &= -\sum_{n \geq 0} \sum_{k \geq 1} \frac{\zeta^k q^{ak+cnk}}{k} = -\sum_{k \geq 1} \frac{\zeta^k q^{ak}}{k \big(1-q^{ck}\big)}.
	\end{align*}
Let $k = mb+j$ to obtain
\begin{align*}
	- \sum_{j=1}^{b} \zeta^j \sum_{m \geq 0} \frac{ q^{a(mb+j)}}{(mb+j) \big(1-q^{c(mb+j)}\big)}.
\end{align*}
We may rewrite this as
\begin{align*}
	- z \sum_{j=1}^{b} \zeta^j \sum_{m \geq 0} f\left( \left(m +\frac{j}{b}\right) bz \right),
\end{align*}
where
\begin{align*}
	f(x) := \frac{{\rm e}^{-ax}}{x \left(1- {\rm e}^{-cx}\right)}.
\end{align*}
Note that $f$ has a Taylor expansion at $x=0$ of $\frac{1}{cx^2} + \frac{\frac{1}{2} - \frac{a}{c}}{x} +O(1)$.
Thus using Lemma \ref{lemma11}, we obtain that
\begin{align*}
		\log F_{a,c}(\zeta;q) = {}& \sum_{j=1}^{b} \zeta^j \biggl( -\frac{1}{b^2 c z} \zeta\left(2,\frac{j}{b} \right) - \frac{I^*_{f,1}}{b}\\
& + \frac{\frac{1}{2} -\frac{a}{c}}{b} \left(\log(bz) + \psi\left(\frac{j}{b}\right) +\gamma \right) +O(|z|) \biggr).
\end{align*}
Then using that $\sum_{j=1}^b \zeta^{j}=0$ and \cite[p.~39]{Campbell} (see \cite{BCMO} which corrects a minus sign and erroneous~$k$ on the right-hand side of Campbell)
\begin{align*}%\label{digam}
	\sum_{j=1}^b \psi \left( \frac{j}{b} \right) \zeta^{j}=b \log ( 1-\zeta )
\end{align*}
along with
\begin{align*}
	\sum_{j =1}^{b} \zeta^j \zeta\left(2 , \frac{j}{b}\right) = b^2 \zeta \Phi(\zeta,2,1)
\end{align*}
and exponentiating gives us that
\begin{align*}
		F_{a,c}(\zeta;q) = (1-\zeta)^{\frac{1}{2} - \frac{a}{c}}\exp\left( - \frac{\zeta \Phi(\zeta,2,1)}{ c z} \right) ( 1+ O(|z|) )
\end{align*}
as desired.
\end{proof}

\begin{thebibliography}{99}
\footnotesize\itemsep=0pt

\bibitem{BCMO}
Bringmann K., Craig W., Males J., Ono K., Distributions on partitions arising
 from {H}ilbert schemes and hook lengths, \href{https://doi.org/10.1017/fms.2022.45}{\textit{Forum Math. Sigma}}
 \textbf{10} (2022), e49, 30~pages, \href{https://arxiv.org/abs/2109.10394}{arXiv:2109.10394}.

\bibitem{Campbell}
Campbell R., Les int\'egrales eul\'eriennes et leurs applications. \'{E}tude
 approfondie de la fonction gamma, \textit{Collection Universitaire de
 Math\'ematiques}, Vol.~20, Dunod, Paris, 1966.

\bibitem{Chern}
Chern S., Nonmodular infinite products and a conjecture of {S}eo and {Y}ee,
 \href{https://doi.org/10.1016/j.aim.2023.108932}{\textit{Adv. Math.}} \textbf{417} (2023), 108932, 40~pages,
 \href{https://arxiv.org/abs/1912.10341}{arXiv:1912.10341}.

\bibitem{Ciolan}
Ciolan A., Equidistribution and inequalities for partitions into powers,
 \textit{Bull. Math. Soc. Sci. Math. Roumanie~(N.S.)} \textbf{66} (2023),
 409--431, \href{https://arxiv.org/abs/2002.05682}{arXiv:2002.05682}.

\bibitem{C}
Craig W., Seaweed algebras and the index statistic for partitions,
 \href{https://doi.org/10.1016/j.jmaa.2023.127544}{\textit{J.~Math. Anal. Appl.}} \textbf{528} (2023), 127544, 23~pages,
 \href{https://arxiv.org/abs/2112.09269}{arXiv:2112.09269}.

\bibitem{HardyRamanujan}
Hardy G.H., Ramanujan S., Asymptotic formulae in combinatory analysis,
 \href{https://doi.org/10.1112/plms/s2-17.1.75}{\textit{Proc. Lond. Math. Soc.}} \textbf{17} (1918), 75--115.

\bibitem{LZ}
Liu Z.-G., Zhou N.H., Uniform asymptotic formulas for the {F}ourier coefficients
 of the inverse of theta functions, \href{https://doi.org/10.1007/s11139-021-00409-8}{\textit{Ramanujan~J.}} \textbf{57} (2022),
 1085--1123, \href{https://arxiv.org/abs/1903.00835}{arXiv:1903.00835}.

\bibitem{NgoRhoades}
Ngo H.T., Rhoades R.C., Integer partitions, probabilities and quantum modular
 forms, \href{https://doi.org/10.1186/s40687-017-0102-4}{\textit{Res. Math. Sci.}} \textbf{4} (2017), 17, 36~pages.

\bibitem{SZ}
Schlosser M.J., Zhou N.H., Expansions of averaged truncations of basic
 hypergeometric series, \href{https://arxiv.org/abs/2307.10821}{arXiv:2307.10821}.

\bibitem{Wright1}
Wright E.M., Stacks, \href{https://doi.org/10.1093/qmath/19.1.313}{\textit{Quart.~J. Math. Oxford Ser.~(2)}} \textbf{19}
 (1968), 313--320.

\bibitem{Wright2}
Wright E.M., Stacks.~{II}, \href{https://doi.org/10.1093/qmath/22.1.107}{\textit{Quart.~J. Math. Oxford Ser. (2)}} \textbf{22}
 (1971), 107--116.

\bibitem{Zhou}
Zhou N.H., Note on partitions into polynomials with number of parts in an
 arithmetic progression, \href{https://doi.org/10.1142/S1793042121500718}{\textit{Int.~J. Number Theory}} \textbf{17} (2021),
 1951--1963, \href{https://arxiv.org/abs/1909.13549}{arXiv:1909.13549}.

\end{thebibliography}
\end{document}
